\documentclass[a4paper,10pt]{article}
\usepackage{subfiles}
\usepackage[margin=1in]{geometry}
\usepackage{adjustbox}
\usepackage{amsfonts}
\usepackage{amsmath}
\usepackage{amssymb}
\usepackage{amsthm}
\usepackage{calc}
\usepackage{enumitem}
\usepackage{eucal}
\usepackage{fancyhdr}
\usepackage{mathrsfs}
\usepackage{mathtools}
\usepackage{nameref}
\usepackage{parskip}
\usepackage{physics}
\usepackage{setspace}
\usepackage{tabularx}
\usepackage[most]{tcolorbox}
\usepackage{tikz}
\usepackage{tikz-cd}
\usepackage{titlesec}
\usepackage{verbatim}
\usepackage{xcolor}
\usepackage{quiver}
\usepackage{imakeidx}
\usepackage{hyperref}


\usetikzlibrary{calc}
\usetikzlibrary{decorations.pathmorphing}
\usetikzlibrary{spath3}

% A TikZ style for curved arrows of a fixed height, due to AndreC.
\tikzset{curve/.style={settings={#1},to path={(\tikztostart)
    .. controls ($(\tikztostart)!\pv{pos}!(\tikztotarget)!\pv{height}!270:(\tikztotarget)$)
    and ($(\tikztostart)!1-\pv{pos}!(\tikztotarget)!\pv{height}!270:(\tikztotarget)$)
    .. (\tikztotarget)\tikztonodes}},
    settings/.code={\tikzset{quiver/.cd,#1}
        \def\pv##1{\pgfkeysvalueof{/tikz/quiver/##1}}},
    quiver/.cd,pos/.initial=0.35,height/.initial=0}

% A TikZ style for shortening paths without the poor behavior of `shorten <' and `shorten >'.
\tikzset{between/.style n args={2}{/tikz/execute at end to={
    \tikzset{spath/split at keep middle={current}{#1}{#2}}
}}}

% TikZ arrowhead/tail styles.
\tikzset{tail reversed/.code={\pgfsetarrowsstart{tikzcd to}}}
\tikzset{2tail/.code={\pgfsetarrowsstart{Implies[reversed]}}}
\tikzset{2tail reversed/.code={\pgfsetarrowsstart{Implies}}}
% TikZ arrow styles.
\tikzset{no body/.style={/tikz/dash pattern=on 0 off 1mm}}

%%%%%%%%%%%%%%%%%%%%%%%%%%%%%%%%%%%%%%%%%%%%%%%%%%%%%%%%%%%%%%%%%%%%%%%%%%%%%%%%

\hypersetup{
    colorlinks,
    citecolor=black,
    filecolor=black,
    linkcolor=blue!60!black,
    urlcolor=blue!60!black
}
\makeindex[title=Index]
\makeindex[name=thm,title=Index of Theorems]


\titleformat{name=\section}
{\normalfont\Large\bfseries}
{}
{0pt}
{}

\tcbuselibrary{theorems}
\NewTcbTheorem[number within=section]{theorem}{Theorem}{
colback=white,
colframe=gray!30,
fonttitle=\bfseries,
coltitle=black,
colbacktitle=gray!30,
toptitle=1mm,
bottomtitle=1mm,
parbox=false}{tcb}
\NewTcbTheorem[use counter from=theorem]{lemma}{Lemma}{
colback=white,
colframe=gray!30,
fonttitle=\bfseries,
coltitle=black,
colbacktitle=gray!30,
toptitle=1mm,
bottomtitle=1mm,
parbox=false}{tcb}
\NewTcbTheorem[use counter from=theorem]{exercise}{Exercise}{
colback=white,
colframe=gray!30,
fonttitle=\bfseries,
coltitle=black,
colbacktitle=gray!30,
toptitle=1mm,
bottomtitle=1mm,
parbox=false}{tcb}
\NewTcbTheorem[use counter from=theorem]{corollary}{Corollary}{
colback=white,
colframe=gray!30,
fonttitle=\bfseries,
coltitle=black,
colbacktitle=gray!30,
toptitle=1mm,
bottomtitle=1mm,
parbox=false}{tcb}
\NewTcbTheorem[use counter from=theorem]{proposition}{Proposition}{
colback=white,
colframe=gray!30,
fonttitle=\bfseries,
coltitle=black,
colbacktitle=gray!30,
toptitle=1mm,
bottomtitle=1mm,
parbox=false}{tcb}

\theoremstyle{definition}
\newtheorem{remark}[tcb@cnt@theorem]{Remark}
\newtheorem{example}[tcb@cnt@theorem]{Example}
\newtheorem{warning}[tcb@cnt@theorem]{Warning}
\newtheorem{notation}[tcb@cnt@theorem]{Notation}
\newtheorem{definition}[tcb@cnt@theorem]{Definition}
\newtheorem{motivation}[tcb@cnt@theorem]{Motivation}
\newtheorem{application}[tcb@cnt@theorem]{Application}
\newtheorem{construction}[tcb@cnt@theorem]{Construction}

\newcommand\aref[2][blue!60!black]{\begingroup\hypersetup{linkcolor=#1}\hyperlink{#2}{#2}\endgroup}
\newcommand\bref[3][blue!60!black]{\begingroup\hypersetup{linkcolor=#1}\hyperlink{#2}{#3}\endgroup}
\newcommand\lref[2][blue!60!black]{lemma \ref{tcb:#2}}
\newcommand\tref[2][blue!60!black]{theorem \ref{tcb:#2}}
\newcommand\eref[2][blue!60!black]{exercise \ref{tcb:#2}}
\newcommand\cref[2][blue!60!black]{corollary \ref{tcb:#2}}
\newcommand\pref[2][blue!60!black]{proposition \ref{tcb:#2}}
\newcommand\Lref[2][blue!60!black]{Lemma \ref{tcb:#2}}
\newcommand\Tref[2][blue!60!black]{Theorem \ref{tcb:#2}}
\newcommand\Eref[2][blue!60!black]{Exercise \ref{tcb:#2}}
\newcommand\Cref[2][blue!60!black]{Corollary \ref{tcb:#2}}
\newcommand\Pref[2][blue!60!black]{Proposition \ref{tcb:#2}}

\makeatletter
\newcommand*{\shortenfences}[3]{%
  % #1: left fence without \left
  % #2: formula inside the fences
  % #3: right fence without \right
  \mathpalette{\@shortenfences{#1}{#3}}{#2}%
}
\newdimen\sf@dimen
\newcommand*{\@shortenfences}[4]{%
  % #1: left fence
  % #2: right fence
  % #3: math style
  % #4: formula
  \sbox0{$#3#4\m@th$}%
  \sbox2{$#3\vcenter{}$}%
  % \dimen0: height above math axis
  \dimen0=\dimexpr\ht0 - \ht2\relax
  \ifdim\dimen0<\z@
    \dimen0=\z@
  \fi
  % \dimen2: depth below math axis
  \dimen2=\dimexpr\ht2 + \dp0\relax
  \ifdim\dimen2<\z@
    \dimen2=\z@
  \fi
  % \sf@dimen: amount for lowering the inner formula
  % to center the inner formula.
  \sf@dimen=\dimexpr(\dimen0 - \dimen2)/2\relax
  % lower the inner formula and raise the outer formula with
  % the fences to keep the math axis.
  \raisebox{\sf@dimen}{%
    $#3\left#1\raisebox{-\sf@dimen}{\box0}\right#2\m@th$%
  }
}
\makeatother

\newcommand\eq{{\operatorname{eq}}}
\newcommand\GL{{\operatorname{GL}}}
\newcommand\Gr{{\operatorname{Gr}}}
\newcommand\id{{\operatorname{id}}}
\newcommand\pr{{\operatorname{pr}}}
\newcommand\pt{{\operatorname{pt}}}
\newcommand\sk{{\operatorname{sk}}}
\newcommand\SL{{\operatorname{SL}}}
\newcommand\cod{{\operatorname{cod}}}
\newcommand\dom{{\operatorname{dom}}}
\newcommand\fin{{\operatorname{fin}}}
\newcommand\Fun{{\operatorname{Fun}}}
\newcommand\Gal{{\operatorname{Gal}}}
\newcommand\Hom{{\operatorname{Hom}}}
\newcommand\lex{{\operatorname{lex}}}
\newcommand\Map{{\operatorname{Map}}}
\newcommand\Mat{{\operatorname{Mat}}}
\newcommand\Mor{{\operatorname{Mor}}}
\newcommand\qis{{\operatorname{qis}}}
\newcommand\red{{\operatorname{red}}}
\newcommand\rev{{\operatorname{rev}}}
\newcommand\cart{{\operatorname{cart}}}
\newcommand\coeq{{\operatorname{coeq}}}
\newcommand\Flag{{\operatorname{Flag}}}
\newcommand\Isom{{\operatorname{Isom}}}
\newcommand\Bilin{{\operatorname{Bilin}}}
\newcommand\canon{{\operatorname{canon}}}
\newcommand\revlex{{\operatorname{revlex}}}


\newcommand\Bl{\operatorname{Bl}}
\newcommand\Cl{\operatorname{Cl}}
\newcommand\gr{\operatorname{gr}}
\newcommand\im{\operatorname{im}}
\newcommand\pd{\operatorname{pd}}
\newcommand\adj{\operatorname{adj}}
\newcommand\Ann{\operatorname{Ann}}
\newcommand\Ass{\operatorname{Ass}}
\newcommand\Aut{\operatorname{Aut}}
\newcommand\Der{\operatorname{Der}}
\newcommand\Div{\operatorname{div}}
\newcommand\DIV{\operatorname{Div}}
\newcommand\Eig{\operatorname{Eig}}
\newcommand\End{\operatorname{End}}
\newcommand\Ext{\operatorname{Ext}}
\newcommand\Fix{\operatorname{Fix}}
\newcommand\kod{\operatorname{kod}}
\newcommand\Kom{\operatorname{Kom}}
\newcommand\lcm{\operatorname{lcm}}
\newcommand\Nil{\operatorname{Nil}}
\newcommand\Orb{\operatorname{Orb}}
\newcommand\Pic{\operatorname{Pic}}
\newcommand\rad{\operatorname{rad}}
\newcommand\rem{\operatorname{rem}}
\newcommand\Syl{\operatorname{Syl}}
\newcommand\Sym{\operatorname{Sym}}
\newcommand\syz{\operatorname{syz}}
\newcommand\Tor{\operatorname{Tor}}
\newcommand\Bord{\operatorname{Bord}}
\newcommand\CaCl{\operatorname{CaCl}}
\newcommand\coim{\operatorname{coim}}
\newcommand\cone{\operatorname{cone}}
\newcommand\Curl{\operatorname{curl}}
\newcommand\Disc{\operatorname{Disc}}
\newcommand\Frac{\operatorname{Frac}}
\newcommand\Grad{\operatorname{grad}}
\newcommand\mult{\operatorname{mult}}
\newcommand\Proj{\operatorname{Proj}}
\newcommand\sign{\operatorname{sign}}
\newcommand\Spec{\operatorname{Spec}}
\newcommand\Stab{\operatorname{Stab}}
\newcommand\Supp{\operatorname{Supp}}
\newcommand\Triv{\operatorname{Triv}}
\newcommand\coker{\operatorname{coker}}
\newcommand\codim{\operatorname{codim}}
\newcommand\colim{\operatorname{colim}}
\newcommand\depth{\operatorname{depth}}
\newcommand\graph{\operatorname{graph}}
\newcommand\MSpec{\operatorname{MSpec}}
\newcommand\Slice{\operatorname{Slice}}
\newcommand\OrBord{\operatorname{OrBord}}
\newcommand\height{\operatorname{height}}

\newcommand\NL{\mathbb{NL}}
\newcommand\trdeg{\operatorname{tr.deg}}
\newcommand\sext{{\operatorname{\mathscr{E}\hspace{-1pt}xt}}}
\newcommand\shom{{\operatorname{\mathscr{H}\hspace{-1.5pt}om}}}
\newcommand\tcap{\mathbin
{\mathmakebox[\widthof{$\cap$}][c]{\top}\mathmakebox[0pt][r]{\cap}}}
\newcommand\pp[2][]{\frac{\partial#1}{\partial#2}}
\newcommand\ppn[3][]{\frac{\partial^{#3}#1}{\partial#2^{#3}}}
\newcommand\ppe[3][]{{\mathmakebox[\widthof{$\displaystyle\frac{\partial#1}
{\partial#2}|$}][r]{\left.\frac{\partial#1}{\partial#2}\right|}}_{#3}}
\newcommand\ppne[4][]{{\mathmakebox[\widthof{$\displaystyle\frac{\partial#1^{#3}}
{\partial#2^{#3}}|$}][r]{\left.\frac{\partial#1}{\partial#2}\right|}}_{#4}}
\newcommand\tpp[2][]{\partial#1/\partial#2}
\newcommand\dpp[2][]{\displaystyle\frac{\partial#1}{\partial#2}}

\renewcommand\op{{\operatorname{op}}}

\DeclareFontFamily{U}{mathx}{}
\DeclareFontShape{U}{mathx}{m}{n}{<-> mathx10}{}
\DeclareSymbolFont{mathx}{U}{mathx}{m}{n}
\DeclareMathAccent{\widecheck}{0}{mathx}{"71}

\DeclareSymbolFont{bbold}{U}{bbold}{m}{n}
\DeclareSymbolFontAlphabet{\mathbbold}{bbold}
\DeclareMathSymbol{\bbomega}{\mathord}{bbold}{"7F}
\DeclareMathSymbol{\bbk}{\mathord}{bbold}{"6B}
\newcommand\ncat[1]{\mathbbold{#1}}

\newcommand\di[1]{\index{#1}\textbf{\textcolor{orange}{#1}}}
\newcommand\dni[1]{\textbf{\textcolor{orange}{#1}}}
\newcommand\cbox[1]{\colorbox{gray!20}{#1}}
\newcommand\dbox[1]{\colorbox{blue!10}{#1}}
\newcommand\ebox[1]{\colorbox{orange!10}{#1}}
\newcommand\cbl[1]{\colorbox{red!10}{#1}}
\newcommand\inj{\mathrel{\hookrightarrow}}
\newcommand\surj{\mathrel{\twoheadrightarrow}}
\newcommand\ol[1]{\overline{#1}}
\newcommand\ul[1]{\underline{#1}}
\newcommand\wh[1]{\widehat{#1}}
\newcommand\wt[1]{\widetilde{#1}}
\newcommand\wc[1]{\widecheck{#1}}
\renewcommand\leq{\leqslant}
\renewcommand\geq{\geqslant}
\renewcommand\mod[1]{\ (\mathrm{mod}\ #1)}
\renewcommand\preceq{\preccurlyeq}
\renewcommand\succeq{\succcurlyeq}
\newcommand\cat[1]{{\textsf{#1}}}
\newcommand\lsub[2]{\vphantom{#2}_{#1}{#2}}
\newcommand\lsup[2]{\vphantom{#2}^{#1}{#2}}

\newcommand\circled[1]{\tikz[baseline=(char.base)]{
\node[shape=circle,draw,inner sep=2pt](char){#1};}}

\newcommand\pagebreaktoggle{}
\newcommand\nonmath[1]{}

\newcommand\sproj[1]{[Stacks Project \href{https://stacks.math.columbia.edu/tag/#1}{#1}]}
\newcommand\krd[1]{[Kerodon \href{https://kerodon.net/tag/#1}{#1}]}
\newcommand\mbook[1]{[Moduli Book \href{https://book.themoduli.space/tag/#1}{#1}]}
\newcommand\mse[1]{[MSE \href{https://math.stackexchange.com/questions/#1}{#1}]}
\newcommand\mo[1]{[MO \href{https://mathoverflow.net/questions/#1}{#1}]}
\newcommand\arxiv[1]{[arXiv \href{https://arxiv.org/abs/math/#1}{#1}]}

\raggedbottom
\pagestyle{fancy}
\lhead{\bref{toc}{ToC}\quad\bref{index}{Index}}
\chead{}
\rhead{}
\cfoot{\thepage}

\setlength{\headheight}{15pt}
\addtolength{\topmargin}{-3pt}
\setlength{\parindent}{16pt}
\setlength{\parskip}{6pt}
\makeatletter
\g@addto@macro\normalsize{%
  \setlength\abovedisplayskip{6pt}
  \setlength\belowdisplayskip{6pt}
  \setlength\abovedisplayshortskip{3pt}
  \setlength\belowdisplayshortskip{3pt}
}
\makeatother
\setlist{listparindent=\parindent,parsep=\parsep,itemsep=0pt,topsep=0pt,partopsep=0pt}
\renewcommand\arraystretch{1}
\renewcommand\baselinestretch{1.2}
\allowdisplaybreaks
\emergencystretch 3em
\begin{document}
\def\numberline#1{}
\hypertarget{toc}{}
\tableofcontents
\pagebreak

\section{Primary Ideal}
\begin{notation}
All rings will be unital and commutative.
\end{notation}

\begin{definition}
Let $R$ be a ring. 
An ideal $\mathfrak{q}\subseteq R$ is a \di{primary ideal} if 
for any $a,b\in R$ with $ab\in\mathfrak{q}$, either $a\in\mathfrak{q}$ 
or $b^n\in\mathfrak{q}$ for some $n\in\mathbb{Z}_{\geq1}$.
Equivalently, $\mathfrak{q}\subseteq R$ is a primary ideal of 
$R/\mathfrak{q}$ is a ring where every zero-divisor is nilpotent.

The radical of a primary ideal is necessarily a prime ideal.
If $\mathfrak{q}\subseteq R$ is a primary ideal with radical $\mathfrak{p}$, 
then we say $\mathfrak{q}$ is a \dni{$\mathfrak{p}$-primary ideal}.
\end{definition}

\begin{remark}
Radical being prime does not implies the ideal is primary.
For example, consider the ideal $I=(x^2,xy)\subseteq k[x,y]$ for some field $k$.
Then $\sqrt{I}=(x)$ is prime, but $I$ is not primary since 
$xy\in I$ but $x\not\in I$ and $y^n\not\in I$ for any $n\in\mathbb{Z}_{\geq1}$.
Another way to see this is that $xy=0$ in $R/I$,
so $y$ is a zero divisor but it is not nilpotent.

However, radical being maximal is a sufficient condition 
for an ideal to be primary:
\end{remark}

\begin{lemma}{}{lem-radical-is-maximal}
Let $R$ be a ring.
If $\mathfrak{q}\subseteq R$ is an ideal such that 
$\mathfrak{m}:=\sqrt{\mathfrak{q}}$ is a maximal ideal,
then $\mathfrak{q}$ is a $\mathfrak{m}$-primary ideal.
\end{lemma}
\begin{proof}
It suffices to show that every zero-divisor in $R/\mathfrak{q}$ is nilpotent.
Prime ideals of $R/\mathfrak{q}$ is in bijection with 
prime ideals of $R$ containing $\mathfrak{q}$.
Since radical of $\mathfrak{q}$ is the intersection of all prime ideals 
containing $\mathfrak{q}$, it follows that $\mathfrak{m}=\sqrt{\mathfrak{q}}$
implies the only prime ideal of $R$ containing $\mathfrak{q}$ is $\mathfrak{m}$.

Thus, $\mathfrak{m}$ is the unique prime ideal of $R/\mathfrak{q}$.
So $R/\mathfrak{q}$ is a local ring with maximal ideal $\mathfrak{m}$
and all zero-divisors in $R/\mathfrak{q}$ is contained in $\mathfrak{m}$
since any element outside $\mathfrak{m}$ is a unit.
Now $\mathfrak{m}=\sqrt{\mathfrak{q}}$ implies 
the some power of any element of $\mathfrak{m}$ is in $\mathfrak{q}$,
which implies some power of any zero-divisor in $R/\mathfrak{q}$ 
is zero, i.e. nilpotent.
\end{proof}

\begin{lemma}{}{lem-of-loc-of-primary-ideal}
Let $R$ be a ring and $S\subseteq R$ be a multiplicative subset.
If $R$ is a ring where every zero-divisor is nilpotent,
then $S^{-1}R$ is also a ring where every zero-divisor is nilpotent.
\end{lemma}
\begin{proof}
Suppose $(a/s)(b/t)=0/1$ in $S^{-1}R$, then $(ab)/(st)=0/1$ 
implies there exists some $u\in S$ such that $u(ab)=0$ in $R$.
So $a$ is a zero divisor in $R$ and 
the assumption implies $a^n=0$ for some $n\in\mathbb{Z}_{\geq1}$.
Then $a/s\in S^{-1}R$ is nilpotent because $(a/s)^n=a^n/s^n=0/1$.
\end{proof}

\begin{theorem}{}{thm-loc-of-primary-ideal}
Let $R$ be a ring and $S\subseteq R$ be a multiplicative subset.
If $\mathfrak{q}\subseteq R$ is a $\mathfrak{p}$-primary ideal 
and $\mathfrak{p}\cap S=\varnothing$,
then $S^{-1}\mathfrak{q}\subseteq S^{-1}R$ 
is a $S^{-1}\mathfrak{p}$-primary ideal.
\end{theorem}
\begin{proof}
\cbox{($S^{-1}\mathfrak{q}$ is a primary ideal)}
It suffices to show $S^{-1}R/S^{-1}\mathfrak{q}$ is a 
ring where every zero-divisor is nilpotent.
Since localization is an exact functor, we have 
$S^{-1}R/S^{-1}\mathfrak{q}\cong S^{-1}(R/\mathfrak{q})$
as an $S^{-1}R$-module.
Both sides are also $S^{-1}R$-algebras and the module isomorphism 
respect multiplication, so this is in fact an isomorphism of rings.
Then by \lref{lem-of-loc-of-primary-ideal}, we have 
\begin{equation*}\begin{split}
\mathfrak{q}\subseteq R\text{ is a primary ideal}&\Longrightarrow
R/\mathfrak{q}\text{ is a ring where every zero-divisor is nilpotent}\\
&\Longrightarrow S^{-1}(R/\mathfrak{q})\text{ is a ring where every zero-divisor is nilpotent}\\
&\Longrightarrow S^{-1}R/S^{-1}\mathfrak{q}\text{ is a ring where every zero-divisor is nilpotent}\\
&\Longrightarrow S^{-1}\mathfrak{q}\subseteq S^{-1}R\text{ is a primary ideal}
\end{split}\end{equation*}
\cbox{($\sqrt{S^{-1}\mathfrak{q}}=S^{-1}\mathfrak{p}$)}
$\sqrt{S^{-1}\mathfrak{q}}$ is the intersection of all prime ideal 
of $S^{-1}R$ containing $S^{-1}\mathfrak{q}$.
So $\sqrt{S^{-1}\mathfrak{q}}\subseteq S^{-1}\mathfrak{p}$ because 
$S^{-1}\mathfrak{p}$ is a prime ideal of $S^{-1}R$ 
containing $S^{-1}\mathfrak{q}$.
This uses the assumption that $S\cap\mathfrak{p}=\varnothing$.

Conversely, suppose $a/s\in S^{-1}\mathfrak{p}$
with $a\in\mathfrak{p}$ and $s\in S$.
Then $a^n\in\mathfrak{q}$ for some $n\in\mathbb{Z}_{\geq1}$
because $\mathfrak{p}=\sqrt{\mathfrak{q}}$.
$S$ is multiplicative and $s\in S$ implies $s^n\in S$.
It follows that $(a/s)^n=a^n/s^n\in S^{-1}\mathfrak{q}$
which shows $a/s\in\sqrt{S^{-1}\mathfrak{q}}$.
This completes the proof because every element of $S^{-1}\mathfrak{p}$
is a finite sum of $a/s$ with $a\in\mathfrak{p}$ and $s\in S$.
\end{proof}

\begin{proposition}{}{prop-loc-of-primary-ideal}
Let $R$ be a ring and $S\subseteq R$ be a multiplicative subset.
Let $\mathfrak{p}$ be a prime ideal of $R$ 
with $\mathfrak{p}\cap S=\varnothing$.
Then every $S^{-1}\mathfrak{p}$-primary ideal 
is of the form $S^{-1}\mathfrak{q}$ where 
$\mathfrak{q}\subseteq R$ is a $\mathfrak{p}$-primary ideal.
\end{proposition}
\begin{proof}
Let $\mathfrak{a}$ be a $S^{-1}\mathfrak{q}$-primary ideal and define 
\begin{equation*}
\mathfrak{q}=\left\{a\in R\;\middle|\;\frac a1\in\mathfrak{a}\right\}\subseteq R
\end{equation*}
Then $\mathfrak{q}\subseteq R$ is an ideal and 
$\mathfrak{a}=S^{-1}\mathfrak{q}$.

To see $\mathfrak{q}$ is a $\mathfrak{p}$-primary ideal,
suppose $ab\in\mathfrak{q}$.
Then $ab/1=(a/1)(b/1)\in S^{-1}\mathfrak{q}=\mathfrak{a}$
which is a $S^{-1}\mathfrak{p}$-primary ideal.
So either $a/1\in\mathfrak{a}$ or $(b/1)^n=b^n/1\in\mathfrak{a}$ for some
$n\in\mathbb{Z}_{\geq1}$.
Therefore, either $a\in\mathfrak{q}$ or $b^n\in\mathfrak{q}$ for some
$n\in\mathbb{Z}_{\geq1}$, which shows $\mathfrak{q}$ is a primary ideal.
Finally, $\mathfrak{q}$ is a $\sqrt{\mathfrak{q}}$-primary ideal 
implies $S^{-1}\mathfrak{q}$ is a $(S^{-1}\sqrt{\mathfrak{q}})$-primary ideal
by \tref{thm-loc-of-primary-ideal}.
This shows 
\begin{equation*}
S^{-1}\sqrt{\mathfrak{q}}=\sqrt{S^{-1}\mathfrak{q}}
=\sqrt{\mathfrak{a}}=S^{-1}\mathfrak{p}
\end{equation*}
Since prime ideal of $S^{-1}R$ is in bijection 
with prime ideal of $R$ disjoint from $S$, we conclude 
$\sqrt{\mathfrak{q}}=\mathfrak{p}$.
\end{proof}

\section{Length of Module}
\begin{definition}
Let $R$ be a ring and $M$ be an $R$-module.
$M$ is a \di{simple module} of $M\not=0$ and the only submodule of 
$M$ is zero and $M$ itself.
A \di{chain} of $M$ is a sequence of submodules 
\begin{equation*}
M=M_0\supsetneq M_1\supsetneq M_2\supsetneq\cdots\supsetneq M_n=0
\end{equation*}
A \di{composition series} is a maximal chain of $M$,
in particular, quotients between consecutive submodules are simple modules.
The \di{length} of $M$ is 
\begin{equation*}
\ell_R(M)=\inf\{n\in\mathbb{Z}_{\geq0}\mid\text{there exist a 
composition series }M=M_0\supsetneq M_1\supsetneq\cdots\supsetneq M_n=0\}
\end{equation*}
\end{definition}

\begin{lemma}{}{lem-submod-has-smaller-length}
Let $R$ be a ring.
If $N\subsetneq M$ is an strict inclusion of $R$-modules.
Then either $\ell_R(N)=\ell_R(M)=\infty$ or 
$\ell_R(N)<\ell_R(M)$.
\end{lemma}
\begin{proof}
Let $M=M_0\supsetneq M_1\supsetneq M_2\supsetneq\cdots\supsetneq M_n=0$ 
be a composition series of $M$ of minimal length.
Let $N_i=N\cap M_i$.
Then 
\begin{equation*}
N=N_0\supseteq N_1\supseteq N_2\supseteq\cdots\supseteq N_n=0
\end{equation*}
Note that $N_i/N_{i+1}=(M_i\cap N)/(M_{i+1}\cap N)\subseteq M_i/M_{i+1}$.
Thus $N_i/N_{i+1}$ is either $0$ or simple. 
If $N_i/N_{i+1}=M_i/M_{i+1}$ for all $i$, then $N=M$.
Therefore, $N_i/N_{i+1}=0$ for some $i$, 
which shows $\ell_R(N)<\ell_R(M)$.
\end{proof}

\begin{theorem}{}{thm-length-well-defined}
Let $R$ be a ring and $M$ be an $R$-module 
with a composition series of length $n$.
Then 
\begin{enumerate}[label=(\arabic*)]
\item
Every composition series of $M$ has length $n$.
\item
Every chain of $M$ has length at most $n$.
\item 
Every chain of $M$ can be extended to a composition series of $M$.
\end{enumerate}
\end{theorem}
\begin{proof}
We may assume $n$ is the smallest length of a composition series of $M$.
Let $M=M_0\supsetneq$$M_1\supsetneq M_2\supsetneq\cdots\supsetneq M_n=0$ 
be the given composition series of $M$.

If $M=M_0'\supsetneq M_1'\supsetneq M_2'\supsetneq\cdots\supsetneq M_m'=0$ 
is another composition series of $M$, then $m\geq n$.
On the other hand, $\ell(M_1')\leq\ell(M)-1$, $\ell(M_2')\leq\ell(M)-2$, 
$\ell(M_3')\leq\ell(M)-3$ and so on by \lref{lem-submod-has-smaller-length}.
Thus, $0\leq\ell(M_m')\leq\ell(M)-m=n-m$ which implies $m\leq n$.
The remaining parts are clear.
\end{proof}

\begin{corollary}{Length is Additive}{cor-length-additive}
Let $R$ be a ring and let $0\to M_1\to M_2\to M_3\to0$ 
be a short exact sequence of $R$-modules.
Then 
\begin{equation*}
\ell_R(M_2)=\ell_R(M_1)+\ell_R(M_3).
\end{equation*}
\end{corollary}

\begin{lemma}{}{lem-simple-module}
Let $R$ be a ring and $M$ be an $R$-module.
The following are equivalent:
\begin{enumerate}[label=(\arabic*)]
\item
$M$ is simple.
\item
$\ell_R(M)=1$.
\item 
$M\cong R/\mathfrak{m}$ for some maximal ideal $\mathfrak{m}\subseteq R$.
\end{enumerate}
\end{lemma}
\begin{proof}
$(3)\Longrightarrow(2)\Longrightarrow(1)$ is clear.
For $(1)\Longrightarrow(3)$, let $x\in M\setminus\{0\}$.
The submodule generated by $x$ is $M$ itself because $M$ is simple.
Consider the $R$-module homomorphism $\varphi:R\to M$ defined by $r\mapsto rx$.
$\varphi$ is surjective because $x$ generate $M$,
and $\ker\varphi$ is a maximal ideal of $R$ because $M$ is simple 
and the image of any ideal containing $\ker\varphi$ would be a submodule of $M$.
\end{proof}

\begin{remark}
In a local ring $R$ with unique maximal ideal $\mathfrak{m}$,
any simple $R$-module is isomorphic to the residue field $R/\mathfrak{m}$.
\end{remark}

\begin{lemma}{}{lem-fg-mod-ann-by-max-ideal}
Let $(R,\mathfrak{m},k)$ be a local ring.
If $M$ is a finitely generated $R$-module annihilated by $\mathfrak{m}$
so that it can be viewed as an $k=R/\mathfrak{m}$-module, 
then $\ell_R(M)=\dim_kM$.
\end{lemma}
\begin{proof}
Let $x_1,\dots,x_n$ be a basis of $M$ as an $k$-module/vector space.
The $x_1,\dots,x_n$ also generate $M$ as an $R$-module because any element 
of $M$ can be written as a $R$-linear combination of $x_1,\dots,x_n$
(write as a $k$-linear combination first, then lift the coefficients to $R$).

For any $i\in\{0,\dots,n\}$, let $M_i$ denote the 
submodule generated by $x_i+1,\dots,x_n$ (for example $M_0=M$, $M_n=0$).
Then $M=M_0\supseteq M_1\supseteq M_2\supseteq\cdots\supseteq M_n=0$ 
is a chain of $M$ and $M_i/M_{i+1}\cong k$ is generated by $x_{i+1}$.
Since $\mathfrak{m}$ annihilates $x_{n+1}$,
it follows that the $R$-module homomorphism 
\begin{equation*}\begin{array}{r@{}c@{}l}
R&{}\longrightarrow{}& M_i/M_{i+1} \\
r&{}\longmapsto{}& rx_{i+1}
\end{array}\end{equation*} 
is surjective with kernel containing $\mathfrak{m}$.
So $M_i/M_{i+1}$ is either zero or isomorphic to $R/\mathfrak{m}=k$.
If $M_i/M_{i+1}=0$, then $x_{i+1}$ is in the $R$-submodule of 
$M$ generated by $x_{i+2},\dots,x_n$,
contradicting the $k$-linear independence of $x_1,\dots,x_n$.

Thus $M_i/M_{i+1}$ is isomorphic to $k$ which implies
they are simple $R$-modules by \lref{lem-simple-module}.
So $M=M_0\supseteq M_1\supseteq M_2\supseteq\cdots\supseteq M_n=0$
is a composition series of $M$ and $\ell_R(M)=n=\dim_kM$
by \tref{thm-length-well-defined}.
\end{proof}

\begin{lemma}{}{lem-finite-length-local-ring}
A noetherian local ring $(R,\mathfrak{m})$ has finite length (as an $R$-module) 
if and only if $\mathfrak{m}^n=0$ for some $n\in\mathbb{Z}_{\geq1}$.
\end{lemma}
\begin{proof}
If $\mathfrak{m}^n\not=0$ for all $n\in\mathbb{Z}_{\geq1}$, 
then $\ell_R(R)=\infty$ because $R\supsetneq\mathfrak{m}
\supsetneq\mathfrak{m}^2\supsetneq\mathfrak{m}^3\supsetneq\cdots$ 
is an infinite descending chain of submodules.
Each inclusion is strict because if $\mathfrak{m}^n=\mathfrak{m}^{n+1}$,
then $\mathfrak{m}^n=0$ by Nakayama's lemma.

Conversely, if $\mathfrak{m}^n=0$ for some $n\in\mathbb{Z}_{\geq1}$,
then consider the chain $R\supsetneq\mathfrak{m}\supsetneq\mathfrak{m}^2
\supsetneq\cdots\supsetneq\mathfrak{m}^{n-1}\supsetneq\mathfrak{m}^n=0$.
Each quotient $\mathfrak{m}^i/\mathfrak{m}^{i+1}$ is finitely generated 
because $R$ is noetherian and is annihilated by $\mathfrak{m}$.
By \lref{lem-fg-mod-ann-by-max-ideal}, $\mathfrak{m}^i/\mathfrak{m}^{i+1}$ 
have finite length finite.
So we can extend the chain to a composition series of $R$ 
with finitely many terms.
Therefore, $R$ has finite length.
\end{proof}

\begin{definition}
Let $(R,\mathfrak{m})$ be a local ring.
A \dni{parameter ideal for $R$}\index{parameter ideal} 
is an ideal $\mathfrak{q}\subseteq\mathfrak{m}$ such that 
$R/\mathfrak{q}$ has finite length.
\end{definition}

\begin{lemma}{}{lem-parameter-ideal}
Let $(R,\mathfrak{m})$ be a noetherian local ring.
An ideal $\mathfrak{q}\subseteq\mathfrak{m}$ is a parameter ideal for $R$
if and only if $\mathfrak{m}^n\subseteq\mathfrak{q}$ for some 
$n\in\mathbb{Z}_{\geq1}$. 
\end{lemma}
\begin{proof}
$R/\mathfrak{q}$ is a local ring with maximal ideal $\mathfrak{m}/\mathfrak{q}$.
So $R/\mathfrak{q}$ has finite length if and only if 
$(\mathfrak{m}/\mathfrak{q})^n=0$ for some $n\in\mathbb{Z}_{\geq1}$
by \lref{lem-finite-length-local-ring}.
Now the ideal $(\mathfrak{m}/\mathfrak{q})^n$ in the ring $R/\mathfrak{q}$
is precisely $\mathfrak{m}^n/\mathfrak{q}$.
Hence the condition that $(\mathfrak{m}/\mathfrak{q})^n=0$
is equivalent to $\mathfrak{m}^n/\mathfrak{q}=0$
which is equivalent to $\mathfrak{m}^n\subseteq\mathfrak{q}$.
\end{proof}

\begin{corollary}{}{cor-parameter-ideal-for-the-ring}
Let $(R,\mathfrak{m})$ be a noetherian local ring.
Then $\mathfrak{q}$ is a parameter ideal for $R$ 
if and only if $\mathfrak{q}$ is a $\mathfrak{m}$-primary ideal.
\end{corollary}
\begin{proof}
$R$ is noetherian implies $\mathfrak{m}$ is finitely generated.
By \lref{lem-parameter-ideal} and \lref{lem-radical-is-maximal}, we have
\begin{equation*}\begin{split}
\mathfrak{q}\text{ is a parameter ideal for }R
&\iff\mathfrak{m}^n\subseteq\mathfrak{q}\text{ for some }n\in\mathbb{Z}_{\geq1}\\
&\iff\sqrt{\mathfrak{q}}=\mathfrak{m}\\
&\iff\mathfrak{q}\text{ is a }\mathfrak{m}\text{-primary ideal}\qedhere
\end{split}\end{equation*}
\end{proof}

\begin{corollary}{}{}
Let $k$ be a field and $R=k[x_1,\dots,x_n]$ be a polynomial ring over $k$.
Let $\mathfrak{m}$ be the maximal ideal 
generated by the variables $x_1,\dots,x_n$.
Then the parameter ideals for the local ring $R_{\mathfrak{m}}$ are precisely
the localizations of $\mathfrak{m}$-primary ideals of $R$.
\end{corollary}
\begin{proof}
Parameter ideals for $R_{\mathfrak{m}}$ are 
$\mathfrak{m}R_{\mathfrak{m}}$-primary ideals of $R_{\mathfrak{m}}$
by \cref{cor-parameter-ideal-for-the-ring}.
Every $\mathfrak{m}R_{\mathfrak{m}}$-primary ideal of $R_{\mathfrak{m}}$ is
the localization of a $\mathfrak{m}$-primary ideal of $R$ 
by \pref{prop-loc-of-primary-ideal}.
\end{proof}

\section{Hilbert-Samuel Polynomial}

Goals:
\begin{enumerate}[label=(\arabic*)]
\item
Prove some result saying length can be computed by dimension over 
some field in some situations.
\item
Prove some result that expresses the Hilbert-Samuel polynomial 
of a module over some local ring in terms of things Macaulay 2 can compute.
\end{enumerate}

\pagebreak

\printindex
\hypertarget{index}{}
\phantomsection
\addcontentsline{toc}{section}{Index}
\end{document}